\documentclass[11pt]{article}

\usepackage[margin=1in]{geometry}
\usepackage{amsmath,amssymb,mathtools}
\usepackage{booktabs,longtable}

\newcommand{\F}{\mathbb F}
\newcommand{\SO}{\operatorname{SO}}
\newcommand{\tr}{\operatorname{tr}}
\setlength{\emergencystretch}{4em}

\title{Corrections and Further Results for\\
  \emph{Rational Rotations of Three Dimensional Euclidean Space}}

\begin{document}
\maketitle

\section{Theorem 31: perpendicular directions and spread}

\subsection{Problem}

Definition 30 evaluates the spread using a non-null \(w\perp v\).  The proof
justifies the existence of \(w\) by nondegeneracy of the plane perpendicular
to a \emph{non-null} \(v\).  This does not cover the identity anchor \(v=0\)
or a nonzero null anchor, including the one exhibited over \(\F_5\) in
Example 27.

\subsection{Suggested repair}

For \(v=(t,r,s)\), consider
\[
 (-r,t,0),\qquad(-s,0,t),\qquad(0,-s,r).
\]
All three vectors are perpendicular to \(v\).  If any has nonzero quadrance,
it supplies the required direction.  If all three are null, then
\[
 Q(-r,t,0)+Q(-s,0,t)-Q(0,-s,r)=2t^2=0.
\]
Since the characteristic is not two, \(t=0\).  The vector \((1,0,0)\) is
then perpendicular to \(v\) and has quadrance one.

\subsection{Recommended paper change}

Replace the nondegenerate-plane sentence with this four-candidate argument.
It is valid for zero, null, and non-null anchors and does not require a basis
choice.

\section{Theorem 40: the missing half-turn classification}

\subsection{Problem}

The proof says that a product of perpendicular reflections is the displayed
half turn by a classification in Section 4.  That section defines the matrix
and proves elementary properties, but contains no converse saying that an
arbitrary rotation of half spread one is a displayed half turn.  The same
paragraph invokes an orthogonal non-null basis without constructing it.

\subsection{Suggested repair}

If \(p(\rho)=1\), the matrix identities used in the paper give
\[
 \tr(\rho)=-1,\qquad
 (\rho-\rho^T)^2=(1+\tr(\rho))(\rho+\rho^T-2I)=0.
\]
A direct calculation for an antisymmetric \(3\)-by-\(3\) matrix shows that
square zero forces \(\rho-\rho^T=0\).  Hence
\[
 P=\frac{I+\rho}{2}
\]
is symmetric and idempotent, with trace one and zero adjugate.  At least one
diagonal entry is nonzero.  A row \(v\) corresponding to such an entry
satisfies
\[
 Q(v)\ne0,\qquad P=\frac{vv^T}{Q(v)}.
\]
It follows that
\[
 \rho=2P-I=\rho_{[v]}.
\]

For a non-null \(v_1\perp v\), put \(v_2=v\times v_1\).  Then
\[
 Q(v_2)=Q(v)Q(v_1)\ne0,
\]
and direct multiplication shows that the two corresponding reflection
matrices satisfy \(\sigma_{v_1}\sigma_{v_2}=\rho\).

\subsection{Recommended paper change}

Insert the projector classification before Theorem 40 and replace the
unsupported basis sentence by the displayed cross-product construction.

\section{Theorem 46: incomplete dependent-case split}

\subsection{Problem}

The written proof treats independent \(v,u\), then identifies every remaining
case with \(u=\lambda v\).  If \(v=0\) and \(u\ne0\), the pair is dependent
but \(u\) is not a multiple of \(v\).  The omitted case is elementary because
\(c(A_0)=I\), but it is not covered by the stated split.

\subsection{Suggested repair}

The case split is unnecessary.  Under
\[
 1+Q(v)\ne0,\qquad 1+Q(u)\ne0,\qquad
 1-v\mathbin{\cdot}u\ne0,
\]
direct denominator-cleared calculation gives
\[
 c(A_v)c(A_u)=c(A_{v\otimes u}),\qquad
 v\otimes u\equiv\frac{v+u-v\times u}{1-v\mathbin{\cdot}u}.
\]
The required domain condition for the output follows from
\[
 1+Q(v\otimes u)=
 \frac{(1+Q(v))(1+Q(u))}
      {(1-v\mathbin{\cdot}u)^2}.
\]
This single calculation covers zero, dependent, independent, null, and
anisotropic anchors.

\subsection{Recommended paper change}

Either add the missing case \(v=0,\ u\ne0\), or replace the case split with
the denominator-cleared calculation above.

\section{Corollary 49: missing cap-domain hypotheses}

\subsection{Problem}

The statement assumes only \(Q(a)xy\ne1\), although its left-hand side
contains both \(c(xA_a)\) and \(c(yA_a)\).  Their necessary domain conditions
are
\[
 1+x^2Q(a)\ne0,\qquad 1+y^2Q(a)\ne0.
\]
These conditions are automatic over an ordered field, but not over every
field in the paper's scope.

For example, over \(\F_5\), take
\[
 a=(1,0,0),\qquad x=2,\qquad y=0.
\]
Then \(Q(a)xy=0\ne1\), while
\[
 1+Q(xa)=1+4=0.
\]
Thus the first cap on the left-hand side is undefined despite the printed
hypothesis being satisfied.

\subsection{Suggested correction}

Assume
\[
 1+x^2Q(a)\ne0,\qquad
 1+y^2Q(a)\ne0,\qquad
 1-Q(a)xy\ne0.
\]
Then the intended conclusion is well-defined:
\[
 c(xA_a)c(yA_a)=c(zA_a),\qquad
 z=\frac{x+y}{1-Q(a)xy}.
\]
The output domain follows from the same denominator identity used in
Theorem 46.

\subsection{Recommended paper change}

Add the two input-cap conditions, or explicitly restrict the corollary to an
ordered field where they are automatic.

\section{Theorem 50: associativity without caps}

\subsection{Problem}

The statement assumes only that the four anchor products occurring in the
two bracketings exist.  Its proof identifies them with
\(c(A_v)c(A_u)c(A_w)\), silently requiring all three input caps to exist.
That requirement does not follow from the hypotheses.

Over \(\F_5\), let
\[
 v=(2,0,0),\qquad u=w=0.
\]
Then \(Q(v)=4=-1\), so \(c(A_v)\) is undefined, while every denominator in
the two anchor-product bracketings is one.  Thus this is a gap in the proof,
not a counterexample to the theorem's statement.

\subsection{Suggested repair}

Work before normalization with scalar-vector pairs:
\[
 (x,v)\otimes(y,u)\equiv
 \bigl(xy-v\mathbin{\cdot}u,\;xu+yv-v\times u\bigr).
\]
Both coordinates of
\[
 ((x,v)\otimes(y,u))\otimes(z,w)
 \quad\text{and}\quad
 (x,v)\otimes((y,u)\otimes(z,w))
\]
agree by the usual dot- and cross-product identities.  This product is
therefore associative without any cap-domain assumption.  Normalizing at
the four partial products appearing in the statement yields Theorem 50 under
exactly its printed hypotheses.

\subsection{Recommended paper change}

Replace the appeal to three input caps by the homogeneous calculation.  The
statement itself need not be strengthened.

\section{Lemma 59: a finite degree-plus-one proof}

Lemma 59 itself does not require an infinite argument.  Suppose that
\(g,h\in\mathbb Q[X_1,\ldots,X_n]\), that \(h\ne0\), and that \(g(a)=0\)
at every rational point \(a\) for which \(h(a)\ne0\).  The paper multiplies
the two polynomials and invokes all rational points.  The following finite
simultaneous-nonvanishing argument proves the same statement.

Regard \(g\) and \(h\) as polynomials in their last variable.  Select one
nonzero coefficient polynomial from each.  Recursively choose values for the
first \(n-1\) variables at which both selected coefficients are nonzero.
After this specialization, the resulting one-variable polynomials \(G,H\)
are nonzero.  Put \(P=GH\) and \(d=\deg P\).  It is enough to test the finite
list
\[
 0,1,\ldots,d.
\]
A nonzero polynomial of degree \(d\) cannot vanish at all \(d+1\) of these
distinct rational values.  Hence one tested value makes \(P\), and therefore
both \(G\) and \(H\), nonzero.  The recursion produces a displayed rational
point at which the original \(g\) and \(h\) are simultaneously nonzero.
This contradicts the hypothesis if \(g\ne0\), so \(g=0\).

This proves the literal statement of Lemma 59 for every supplied finite
number of variables.  It uses only coefficient scans, specialization, and the
finite root bound.  Theorem 58 may still be proved more directly by checking
its denominator-cleared coordinate identities, so it does not need the
transfer lemma as an intermediate step.

\subsection{Recommended paper change}

Replace the appeal to the infinitude of \(\mathbb Q\) with the recursive
degree-plus-one argument above.  The statement of Lemma 59 need not change.

\clearpage
\section{Finite formulations of the size claims}

Theorem 62 describes the positive rational square classes as an infinite
group, while Corollary 63 describes the unit-representative subgroup as
having infinite index.  For formal verification we isolate the uniform
finite content of these statements.  Given any positive integer \(N\), we
construct \(N\) distinct square classes and \(N\) distinct kernel cosets.
The integer \(N\) is arbitrary, but it is fixed before the calculation
begins; every object used in the proof is then finite and explicit.  No
completed collection, global enumeration, or choice principle is required.

\medskip
\noindent\textbf{Theorem (Finite form of Theorem 62).}
For every positive integer \(N\), there are positive rationals
\(c_1,\ldots,c_N\) and rational rotations
\(\rho_1,\ldots,\rho_N\) such that
\[
 \chi(\rho_i)=[c_i]
 \quad\text{and}\quad
 [c_i]\ne[c_j]\quad\text{whenever }i\ne j.
\]
Moreover, a non-null proportion producing each \(\rho_i\) can be constructed
explicitly.

\smallskip
\noindent\textbf{Proof.}
Construct distinct primes \(p_1,\ldots,p_N\) by finite Euclid steps.  After
\(p_1,\ldots,p_k\) have been found, put
\[
 m=1+p_1\cdots p_k
\]
and take \(p_{k+1}\) to be the least divisor of \(m\) greater than one.
Bounded trial division finds this divisor, and its minimality makes it prime.
It is different from every preceding prime.  Distinct primes represent
distinct rational square classes: if \(p_i/p_j=(a/b)^2\), then
\(p_i b^2=p_j a^2\), contrary to the parity of the exponent of \(p_i\) in
the two sides.

For each \(i\), the constructive four-square theorem gives integers
\(x_i,t_i,r_i,s_i\) satisfying
\[
 x_i^2+t_i^2+r_i^2+s_i^2=p_i.
\]
Thus the proportion \((x_i:t_i,r_i,s_i)\) is non-null, has quadrance \(p_i\),
and produces a rotation \(\rho_i\) satisfying
\(\chi(\rho_i)=[p_i]\).  Taking \(c_i=p_i\) proves the result.
\hfill\(\blacksquare\)

\medskip
\noindent\textbf{Corollary (Finite form of Corollary 63).}
Let \(K=\ker\chi\), equivalently the subgroup of rotations admitting a
rational proportion representative of quadrance one.  For every positive
integer \(N\), there are rational rotations
\(\rho_1,\ldots,\rho_N\) lying in \(N\) distinct cosets of \(K\).

\smallskip
\noindent\textbf{Proof.}
Take the rotations supplied by the theorem.  If \(i\ne j\), then
\[
 \chi(\rho_i^{-1}\rho_j)=[c_i]^{-1}[c_j]\ne[1].
\]
By the obstruction criterion in Corollary 63, the relative rotation
\(\rho_i^{-1}\rho_j\) has no rational proportion representative of quadrance
one.  Thus it does not belong to \(K\), and \(\rho_i\) and \(\rho_j\) lie in
distinct cosets.
\hfill\(\blacksquare\)

These are the forms used in the formal development.  They retain an
arbitrarily large finite challenge while requiring only \(N\) Euclid steps
and \(N\) finite four-square calculations.  In particular, the proof never
forms a completed family of all primes, all square classes, or all kernel
cosets.  This uniform finite construction is how the apparently infinite
conclusions are formalised without any non-constructive argument.

\clearpage
\section{Conjecture 52 and Problem 68}

\subsection{Counterexample}

The printed family is
\[
 v\otimes u=
 \frac{\alpha v+\beta u+\gamma(v\times u)}
      {1+\delta(v\mathbin{\cdot}u)}.
\]
Take
\[
 \alpha=1,\qquad\beta=\gamma=\delta=0.
\]
Then
\[
 v\otimes u=v
\]
for all \(v,u\).  This is a total associative operation.  It cannot be a
nonzero coefficient rescaling of either product proposed in the conjecture,
because its coefficient of \(u\) is zero whereas theirs is nonzero.  Taking
all four coefficients to be zero gives the constant-zero operation as a
second counterexample.

Consequently Conjecture 52, and hence Problem 68 as printed, has a negative
answer.

\subsection{A corrected classification}

The counterexamples show that associativity alone does not distinguish the
anchor products.  Retaining the homogeneous denominator constant and
requiring the zero vector to be an identity gives an exact replacement for
the printed conjecture.

\medskip
\noindent\textbf{Theorem (Unital associative coefficient products).}
Let \(\alpha,\beta,\gamma,\epsilon,\delta\in\mathbb Q\), and define
\begin{equation}
 v\otimes u=
 \frac{\alpha v+\beta u+\gamma(v\times u)}
      {\epsilon+\delta(v\mathbin{\cdot}u)}
\label{eq:unital-coefficient-product}
\end{equation}
whenever the denominator is nonzero.  The following conditions are
equivalent:
\begin{enumerate}
\item[(i)] the operation defined by
\eqref{eq:unital-coefficient-product} has the zero vector
\(0=(0,0,0)\in\mathbb Q^3\) as a two-sided identity, and is associative
whenever all four products occurring in the two bracketings are defined;
\item[(ii)]
\[
 \epsilon\ne0,\qquad \alpha=\beta=\epsilon,
 \qquad \epsilon\delta=-\gamma^2.
\]
\end{enumerate}

\smallskip
\noindent\textbf{Proof.}
Suppose first that condition \emph{(i)} holds.  Substituting the zero vector into
\eqref{eq:unital-coefficient-product}, the identities
\(v\otimes0=v\) and \(0\otimes u=u\) force
\(\epsilon\ne0\) and \(\alpha=\beta=\epsilon\).  Comparing one
coordinate in the associativity equation
\[
 (xe_1\otimes xe_1)\otimes e_2
   =xe_1\otimes(xe_1\otimes e_2)
\]
gives
\[
 (\epsilon\delta+\gamma^2)x^2=0.
\]
This calculation is valid whenever
\(\epsilon+\delta x^2\ne0\).  At least one of \(x=1\) and \(x=2\)
satisfies this condition: otherwise \(\epsilon+\delta=0\) and
\(\epsilon+4\delta=0\), which would imply \(\epsilon=0\).  Taking
such a nonzero \(x\) gives \(\epsilon\delta=-\gamma^2\).

Conversely, assume the displayed coefficient conditions and put
\(\lambda=\gamma/\epsilon\).  Dividing numerator and denominator by
\(\epsilon\) gives
\[
 v\otimes u=
 \frac{v+u+\lambda(v\times u)}
      {1-\lambda^2(v\mathbin{\cdot}u)}.
\]
The zero vector is visibly a two-sided identity.  If \(\lambda=0\), this is
ordinary vector addition.  If \(\lambda\ne0\), let
\(D_\lambda(v)=\lambda v\).  Direct calculation gives
\[
 D_\lambda(v\otimes u)=
 \frac{D_\lambda(v)+D_\lambda(u)
       +D_\lambda(v)\times D_\lambda(u)}
      {1-D_\lambda(v)\mathbin{\cdot}D_\lambda(u)}.
\]
Thus \(D_\lambda\) transports the product, including its domain of
definition, to the opposite anchor product.  Associativity follows from that
of the opposite anchor product and the invertibility of \(D_\lambda\).
\hfill\(\blacksquare\)

\medskip
\noindent\textbf{Corollary (Normalized forms).}
After normalizing \(\epsilon=1\), every unital associative member of the
family has
\[
 \alpha=\beta=1,\qquad \delta=-\gamma^2.
\]
The case \(\gamma=0\) is vector addition.  Every case \(\gamma\ne0\) is an
invertible coordinate dilation of the opposite anchor product.  If the cross
coefficient is further normalized to \(\gamma=\pm1\), the two resulting
operations are exactly the anchor product and its opposite.

\smallskip
\noindent\textbf{Proof.}
Multiplying all five coefficients by the same nonzero scalar does not change
the partial operation.  Scaling by \(\epsilon^{-1}\) gives the displayed
normalization, and the theorem gives \(\delta=-\gamma^2\).  The final claims
follow from the dilation in the proof, with \(\gamma=-1\) giving the anchor
product and \(\gamma=1\) its opposite. \hfill\(\blacksquare\)

\section{Further results prompted by the open problems}

This section records only results actually established.  It does not list the
parts of the open-ended questions for which no conclusion is claimed.

\subsection{Problem 69: finite odd fields}

Let \(F=\F_q\) be a finite field of odd order \(q\).  The counts requested in
Problem 69 can be obtained from finite square tables and the projective
rotation theorem, without appealing to a group-order formula.  We give the
elementary counting statements first.

We use the classical result that \(-1\) is a square in \(F\) precisely when
\(q\equiv1\pmod4\), and is a nonsquare precisely when
\(q\equiv3\pmod4\).

\medskip
\noindent\textbf{Lemma (Binary circle count).}
For \(b\in F\), let
\[
 N(b)=\#\{(r,s)\in F^2:r^2+s^2=b\}.
\]
If \(-1\) is a square, then \(N(0)=2q-1\) and \(N(b)=q-1\) for
\(b\ne0\).  If \(-1\) is a nonsquare, then \(N(0)=1\) and \(N(b)=q+1\)
for \(b\ne0\).

\smallskip
\noindent\textbf{Proof.}
If \(-1=i^2\), the invertible substitution
\((r,s)\mapsto(r+is,r-is)\) identifies the equation with \(xy=b\).  There
are \(2q-1\) pairs with \(xy=0\), and \(q-1\) with \(xy=b\ne0\).

Suppose instead that \(-1\) is a nonsquare.  The equation at \(b=0\) has
only the solution \((0,0)\).  For \(b\ne0\), the two sets
\[
 \{r^2:r\in F\},\qquad \{b-s^2:s\in F\}
\]
both have \((q+1)/2\) elements, and hence intersect.  Thus the projective
conic \(r^2+s^2=bt^2\) has a point.  The lines through that point
parametrize its points by the projective line, giving \(q+1\) points.  None
lies at infinity, since \(r^2+s^2=0\) has no nonzero solution.  Therefore
\(N(b)=q+1\). \hfill\(\blacksquare\)

\medskip
\noindent\textbf{Lemma (Ternary quadric count).}
For \(a\in F\),
\[
 \#\{v\in F^3:Q(v)=a\}=
 \begin{cases}
 q^2,&a=0,\\
 q^2+q,&-a\text{ is a nonzero square in }F,\\
 q^2-q,&-a\text{ is a nonsquare in }F.
 \end{cases}
\]

\smallskip
\noindent\textbf{Proof.}
Write \(\epsilon=1\) when \(-1\) is a square and \(\epsilon=-1\) otherwise.
The binary-circle lemma says
\[
 N(0)=q+(q-1)\epsilon,
 \qquad N(b)=q-\epsilon\quad(b\ne0).
\]
The required number is \(\sum_{t\in F}N(a-t^2)\).  If \(a=0\), exactly one
summand has zero argument, and the sum is \(q^2\).  If \(a\ne0\) is a square,
exactly two summands have zero argument, and the sum is \(q^2+q\epsilon\).
If \(a\) is a nonsquare, no summand has zero argument, and the sum is
\(q^2-q\epsilon\).  In the last two cases the sign is positive exactly when
\(-a\) is a square. \hfill\(\blacksquare\)

\medskip
\noindent\textbf{Theorem (Rotations and half-spread fibers).}
The rotation group and its half-spread fibers have orders
\[
 |\SO(3,\F_q)|=q(q^2-1),
\]
\[
 \#\{\rho:p(\rho)=0\}=q^2,
 \qquad
 \#\{\rho:p(\rho)=1\}=q^2,
\]
and, for \(h\ne0,1\), on putting \(D=-h(1-h)\),
\[
 \#\{\rho:p(\rho)=h\}=
 \begin{cases}
 q^2+q,&D\text{ is a square in }F,\\
 q^2-q,&D\text{ is a nonsquare in }F.
 \end{cases}
\]

\smallskip
\noindent\textbf{Proof.}
By Theorem 58, rotations are the non-null proportions \((x:v)\).  In the
affine chart \(x=1\), the excluded points satisfy \(Q(v)=-1\), and the
preceding lemma counts \(q^2+q\) of them.  At infinity, the non-null vector
lines number \(q^2\): among the \(q^3-1\) nonzero vectors, exactly \(q^2-1\)
are null, and each line has \(q-1\) nonzero representatives.  Therefore
\[
 |\SO(3,\F_q)|=\bigl(q^3-(q^2+q)\bigr)+q^2=q^3-q.
\]

If \(h\ne1\), every rotation in the fiber has the unique affine form
\(\rho(1:v)\), and the half-spread identity gives
\[
 p(\rho(1:v))=h
 \quad\Longleftrightarrow\quad
 Q(v)=\frac{h}{1-h}.
\]
For \(h=0\) the ternary quadric lemma gives \(q^2\).  For \(h\ne0,1\), the
square class of \(-h/(1-h)\) is that of \(-h(1-h)=D\), giving the two middle
fiber counts.  Finally, \(p=1\) consists exactly of the \(q^2\) non-null
proportions at infinity. \hfill\(\blacksquare\)

\medskip
\noindent\textbf{Theorem (Parabolic rotations).}
There are \(q^2-1\) nonidentity parabolic rotations.  They form one conjugacy
class in \(\SO(3,\F_q)\), and a rotation belongs to the zero half-spread
fiber precisely when its characteristic polynomial is \((X-1)^3\).

\smallskip
\noindent\textbf{Proof.}
The affine description of the zero fiber is \(\rho(1:v)\) with \(Q(v)=0\).
The vector \(v=0\) gives the identity, leaving \(q^2-1\) nonidentity
parabolics.  Equivalently, the nonzero null vectors split into \(q+1\)
projective directions with \(q-1\) nonzero scalar parameters over each.
The characteristic polynomial from Theorem 29 becomes \((X-1)^3\) exactly
when \(p=0\).

For conjugacy, take two nonidentity parabolics and recover their unique
nonzero affine null anchors.  The two anchors are joined by a product of two
non-null reflections.  Conjugation by the resulting rotation transports the
first anchor to the second, and equivariance of the cap then transports the
first parabolic to the second. \hfill\(\blacksquare\)

\medskip
\noindent\textbf{Corollary (Excluded quadrics and quadrance classes).}
The excluded affine norm quadric has \(q^2+q\) points and its projective
closure has \((q+1)^2\) points.  The quadrance-class homomorphism is onto the
two square classes, and each fiber has size
\[
 \frac{q(q^2-1)}2.
\]

\smallskip
\noindent\textbf{Proof.}
The affine count is the ternary quadric lemma at \(a=-1\).  The points at
infinity are the \(q+1\) projective null directions, so the projective count
is \(q^2+q+(q+1)=(q+1)^2\).  Every element of \(F\) is a sum of two squares,
which supplies a non-null proportion in either prescribed square class.
Multiplication by a fixed proportion of nonsquare quadrance exchanges the two
fibers bijectively.  They therefore have equal size, and the rotation-count
theorem gives the displayed value. \hfill\(\blacksquare\)

\subsection{Problem 70: the unit-representative kernel}

Problem 70 asks which subgroups of \(SO(3,\mathbb Q)\) are visible through
\(\chi\), asks for a suitable sense in which its kernel is dense, and asks
for both a description of the finitely generated subgroups of the kernel and
the values of \(\chi\) on rotations of finite order.  The following results
give the parts of that programme which we have obtained.

\medskip
\noindent\textbf{Theorem (The unit-representative kernel).}
Let \(K\) be the rational rotations admitting a rational proportion \(P\)
with \(N(P)=1\).  Then
\[
 K=\ker\chi .
\]
In particular, the definition of \(K\) is independent of the representative,
and \(K\) is closed under products and inverses.

\smallskip
\noindent\textbf{Proof.}
Replacing \(P\) by \(\lambda P\), for
\(\lambda\in\mathbb Q^\times\), replaces its quadrance by
\(\lambda^2N(P)\).  Hence the projective class of \(P\) has a unit
representative exactly when \([N(P)]=[1]\), which is exactly the condition
\(\chi(\rho)=[1]\) for the rotation \(\rho\) represented by \(P\).
Multiplicativity of the quadrance gives
\(N(P\otimes R)=N(P)N(R)\), while conjugation preserves the quadrance and
represents the inverse rotation.  The asserted closure follows.
\hfill\(\blacksquare\)

\medskip
\noindent\textbf{Corollary (Square-axis half turns).}
Every half turn about an axis vector \(a\in\mathbb Q^3\) with nonzero square
quadrance belongs to \(K\).

\smallskip
\noindent\textbf{Proof.}
The half turn is represented by \((0:a)\), of quadrance \(Q(a)\).  If
\(Q(a)=r^2\ne0\), then \(r^{-1}(0:a)\) has quadrance one.
\hfill\(\blacksquare\)

\medskip
\noindent\textbf{Lemma (Finite signatures of positive square classes).}
Write a positive rational in reduced form as \(c=m/n\), where
\(m,n\in\mathbb Z_{>0}\).  Let \(S(c)\) be the finite set of primes occurring
to odd exponent in \(mn\), and let
\(\kappa(c)\) be their product.  Writing \(A\mathbin\triangle B\) for the
symmetric difference \((A\setminus B)\cup(B\setminus A)\), we have
\[
 [c]=[\kappa(c)],\qquad
 [c]=[d]\ \Longleftrightarrow\ S(c)=S(d),\qquad
 S(cd)=S(c)\mathbin\triangle S(d).
\]

\smallskip
\noindent\textbf{Proof.}
Finite prime factorization gives \(mn=\kappa(c)s^2\) for an integer \(s\),
and hence
\[
 c=\kappa(c)\left(\frac{s}{n}\right)^2.
\]
Equality of positive rational square classes is therefore equality of the
parities of all prime exponents.  Under multiplication those parities add
modulo two, which is symmetric difference of the two finite supports.
\hfill\(\blacksquare\)

\medskip
\noindent\textbf{Theorem (Finite-word classification through \(\chi\)).}
Let \(g_1,\ldots,g_m\) be supplied rational rotations.  There is a computable
finite prime set \(S\), a vector \(s_i\in\mathbb F_2^S\) for each generator,
and a finite-dimensional span
\[
 V=\langle s_1,\ldots,s_m\rangle\subseteq\mathbb F_2^S
\]
with the following properties.  Every supplied finite signed word \(w\) has
an exactly computable rational rotation matrix and a vector
\(\nu(w)\in V\); moreover
\[
\begin{aligned}
 w\text{ represents an element of }K
   &\Longleftrightarrow \nu(w)=0,\\
 w_1\text{ and }w_2\text{ represent the same }K\text{-coset}
   &\Longleftrightarrow \nu(w_1)=\nu(w_2).
\end{aligned}
\]
If \(r\) is the rank computed by finite row reduction, then \(V\) has exactly
\(2^r\) vectors, and an explicit finite word realizing each vector is
produced.

\smallskip
\noindent\textbf{Proof.}
Choose one rational proportion for each \(g_i\), take the union of the finite
prime supports in the preceding lemma, and use its parity vectors \(s_i\).
Evaluation of a word is finite proportion multiplication.  Concatenation
gives addition in \(\mathbb F_2^S\), and reversing the word and conjugating
its letters gives the inverse without changing its square-class vector.
The preceding theorem identifies the zero vector with \(K\), and applying it
to \(w_1w_2^{-1}\) gives the coset test.  Finite binary row reduction computes
a basis of \(V\), proves the count \(2^r\), and retains a Boolean generator
word for each vector.  Exact comparison of the nine rational matrix entries
separately decides whether two words represent the same rotation.
\hfill\(\blacksquare\)

\medskip
\noindent\textbf{Theorem (Finite-precision approximation by \(K\)).}
For every rational rotation \(\rho\) and positive rational
\(\varepsilon\), one can construct a rational rotation \(\eta\in K\) such
that
\[
 \sum_{i,j=1}^3(\eta_{ij}-\rho_{ij})^2<\varepsilon.
\]

\smallskip
\noindent\textbf{Proof.}
Choose an oriented rational proportion \(P=(x:v)\) for \(\rho\), put
\(n=N(P)>0\), and compute a sufficiently narrow finite rational
square-root bracket for \(n\), with upper endpoint \(u\).  Set
\[
 A=u+x,\qquad
 R=(A^2-Q(v):2Av),\qquad
 D=A^2+Q(v),\qquad U=D^{-1}R.
\]
Direct expansion gives \(N(R)=D^2\), so \(N(U)=1\) and the rotation
\(\eta\) represented by \(U\) lies in \(K\).  The exact error identity is
\[
 \sum_{i,j=1}^3(\eta_{ij}-\rho_{ij})^2
 =\frac{8(u^2-n)^2Q(v)}{D^2n}
 =8p(\eta^{-1}\rho).
\]
The bracket width is chosen by finite integer square-root refinement so that
the displayed rational quantity is below \(\varepsilon\).  This constructs
one certified answer for each finite precision request; it uses no real
square root, sequence, limit, or completed topology.
\hfill\(\blacksquare\)

\medskip
\noindent\textbf{Theorem (Finite-order values of \(\chi\)).}
Suppose that a rational rotation \(\rho\) is supplied with an exact positive
order \(n\), meaning that \(\rho^n=I\) and
\(\rho^j\ne I\) for \(0<j<n\).  Then \(n\) and the corresponding half spread
and quadrance class occur in the following table:
\[
\begin{array}{c|c|c}
\text{exact order}&\text{half spread}&\chi(\rho)\\ \hline
1&0&[1]\\
2&1&[Q(a)]\text{, where }a\text{ is the axis}\\
3&3/4&[1]\\
4&1/2&[2]\\
6&1/4&[3].
\end{array}
\]
In particular, no other exact positive order occurs.

\smallskip
\noindent\textbf{Proof.}
Put \(\tau=\tr(\rho)-1\) and define
\[
 L_0=2,\qquad L_1=\tau,\qquad
 L_{k+2}=\tau L_{k+1}-L_k.
\]
The rational Cayley--Hamilton calculation gives
\(\tr(\rho^k)=1+L_k\), so \(\rho^n=I\) gives \(L_n=2\).  Since
\(L_n-2\) is monic in \(\tau\), clearing the denominator of the reduced
rational \(\tau\) and applying B\'ezout forces that denominator to be one.
The direct half-spread bounds give \(-2\leq\tau\leq2\).  Testing the five
integer values and using the supplied minimality of \(n\) gives respectively
the orders \(2,3,4,6,1\).  Finally
\[
 p=\frac{2-\tau}{4}
\]
and the proportion formula for \(\chi\) give the displayed classes; when
\(p=1\), the representative is \((0:a)\), so its class is
\([Q(a)]\). \hfill\(\blacksquare\)

\medskip
\noindent\textbf{What has not been proved from Problem 70.}
These results do not amount to the full subgroup classification requested in
the problem.  In particular:
\begin{enumerate}
\item We have not classified all subgroups of \(SO(3,\mathbb Q)\) that might
be meant by \emph{visible through \(\chi\)}, nor the full lattice of inverse
images or intersections associated with \(\chi\).
\item For a supplied finite generator list, we classify every supplied finite
word and the complete finite \(\chi\)-image.  We have not classified the
resulting subgroup internally: for example, we do not give its isomorphism
type, defining relations, or a finite presentation.  Thus the request to
describe its finitely generated subgroups remains open beyond the exact
finite-query and quotient result above.
\item The paper leaves \emph{dense in a suitable sense} unspecified.  We prove
the stated finite squared-entrywise (equivalently, squared Frobenius)
approximation theorem, but have not formalized a topological closure statement
or identified this metric formulation with the particular classical topology
intended by the author.
\item The finite-order table determines \(\chi\) completely in orders
\(1,3,4,\) and \(6\).  In order \(2\) it reduces the answer to the square
class \([Q(a)]\) of the rational axis.  We have not classified globally which
positive rational square classes occur as quadrances of rational axes; that
would require the corresponding full rational three-square representability
classification.
\end{enumerate}

\clearpage
\subsection{Problem 71: arbitrary ternary forms}

Problem 71 asks for the theory with the Euclidean quadrance replaced by a
general nondegenerate ternary form, with particular attention to the rational
Minkowski form and its null and parabolic phenomena.  The determinant of the
displayed form cannot in general be omitted from the Euclidean formulas.

\medskip
\noindent\textbf{Theorem (Projective Clifford correspondence for a ternary
form).}
Let \(F\) be a field of characteristic different from two with decidable
equality.  Let \(M\) be a symmetric nonsingular \(3\)-by-\(3\) matrix over
\(F\), and put
\[
 B_M(u,v)=uMv^T,\qquad Q_M(u)=B_M(u,u),\qquad d=\det M,
\]
\[
 C_M(u,v)=(u\mathbin\times v)\operatorname{adj}(M).
\]
For pairs \(P=(x:u)\), define
\[
\begin{split}
 (x:u)\otimes_M(y:v)
   &=
 \bigl(xy-dB_M(u,v):xv+yu-C_M(u,v)\bigr),\\
 N_M(x:u)&=x^2+dQ_M(u),\qquad
 \overline{(x:u)}=(x:-u).
\end{split}
\]
Then \(\otimes_M\) is associative,
\[
 N_M(P\otimes_M R)=N_M(P)N_M(R),
 \qquad
 P\otimes_M\bar P=\bar P\otimes_M P=(N_M(P):0).
\]
For \(N_M(P)\ne0\), the action
\[
 (0:w)\longmapsto
 N_M(P)^{-1}\bigl(\bar P\otimes_M(0:w)\otimes_M P\bigr)
\]
is a member of \(SO(M,F)\).  Every member of \(SO(M,F)\) is obtained in this
way, and two non-null pairs give the same member exactly when one is a
nonzero scalar multiple of the other.

\smallskip
\noindent\textbf{Proof.}
The adjugate identities for \(M\) give the determinant-weighted
Binet--Cauchy and Lagrange identities for \(C_M\), together with the two
vector triple-product identities.  Substitution of those finite coordinate
identities into the displayed product proves associativity and norm
multiplicativity; substituting \(v=-u\) gives the conjugation identities.
The same calculation shows that the normalized conjugation action preserves
\(B_M\) and has determinant one.  Rescaling \(P\) by
\(\lambda\in F^\times\) multiplies both its action numerator and its norm by
\(\lambda^2\), so the action is projectively well defined.

For surjectivity, first select a non-null vector by testing the three
coordinate vectors and their three pairwise sums.  The identity is the
product of the half turn about this vector with itself.  Given a nonidentity
\(\rho\), the same finite coordinate tests produce a non-null \(v\) for
which \(v\rho-v\) is non-null.  Reflection in that displacement reduces
\(\rho\) to an orientation-reversing isometry fixing \(v\).  A finite
selection in \(v^\perp\), followed by reconstruction in an explicit
orthogonal frame, proves that the reduced isometry is one further non-null
reflection.  Thus \(\rho\) is a product of two non-null reflections, hence
also of the corresponding two half turns.  Multiplying their pure pairs
constructs \(P\).

Finally, if two pairs have the same action, multiply one by the conjugate of
the other.  The resulting pair acts trivially.  In characteristic different
from two, the coordinate action formula forces its vector coordinate to
vanish, so it is projectively scalar.  Cancelling the non-null factor gives
the required nonzero scale between the original pairs.
\hfill\(\blacksquare\)

\medskip
\noindent\textbf{Corollary (The form-relative Cayley chart).}
For \(u\in F^3\) with \(1+dQ_M(u)\ne0\), the pair \((1:u)\) defines a
rotation \(c_M(u)\) whose half spread is
\[
 p(u)=\frac{dQ_M(u)}{1+dQ_M(u)}.
\]
If \(u,v\) also satisfy
\[
 1+dQ_M(v)\ne0,\qquad 1-dB_M(u,v)\ne0,
\]
then
\[
\begin{split}
 u\otimes_M v
   &=\frac{u+v-C_M(u,v)}{1-dB_M(u,v)},\\
 c_M(u)c_M(v)&=c_M(u\otimes_M v),
\end{split}
\]
and
\[
 1+dQ_M(u\otimes_M v)
 =
 \frac{(1+dQ_M(u))(1+dQ_M(v))}
      {(1-dB_M(u,v))^2}.
\]
For a rotation of half spread \(p\), the spread between any non-null
direction perpendicular to its recovered projective axis and its image is
\(4p(1-p)\).

\smallskip
\noindent\textbf{Proof.}
The norm of \((1:u)\) is \(1+dQ_M(u)\), and the trace of its normalized
conjugation action gives the displayed half-spread formula.  Multiplying
\((1:u)\) and \((1:v)\) gives scalar coordinate
\(1-dB_M(u,v)\) and vector coordinate
\(u+v-C_M(u,v)\).  Dividing by the nonzero scalar proves the product formula
and the cap-composition identity.  Norm multiplicativity, followed by the
same division, gives the denominator identity.  The final spread formula is
the direct bilinear evaluation of the conjugation action on a perpendicular
non-null direction.
\hfill\(\blacksquare\)

\medskip
\noindent\textbf{Theorem (The rational split model and its null conic).}
For the rational Minkowski form
\[
 Q_M(t,r,s)=t^2+r^2-s^2,
\]
the map
\[
 \Phi(x:t,r,s)=
 \begin{pmatrix}
 x+t&r-s\\
 r+s&x-t
 \end{pmatrix}
\]
is a multiplicative bijection from the even Clifford algebra to
\(\operatorname{Mat}_{2}(\mathbb Q)\), and
\[
 \det\Phi(P)=N_M(P).
\]
Moreover, every nonzero rational null direction has, up to a displayed
nonzero scale, the form
\[
 u(\alpha,\beta)=
 (\alpha^2-\beta^2,\,2\alpha\beta,\,\alpha^2+\beta^2),
 \qquad(\alpha,\beta)\ne(0,0).
\]

\smallskip
\noindent\textbf{Proof.}
The inverse of \(\Phi\) is the finite coordinate map
\[
 \begin{pmatrix}a&b\\c&e\end{pmatrix}
 \longmapsto
 \left(\frac{a+e}{2}:
   \frac{a-e}{2},\,\frac{b+c}{2},\,\frac{c-b}{2}\right).
\]
Direct multiplication proves that \(\Phi(P\otimes_M R)=\Phi(P)\Phi(R)\),
and expansion of the \(2\)-by-\(2\) determinant gives \(N_M(P)\).

The identity \(Q_M(u(\alpha,\beta))=0\) is immediate by expansion.  Conversely,
let \(w=(t,r,s)\ne0\) satisfy \(Q_M(w)=0\).  If \(s+t=0\), then the null
equation gives \(r=0\), and \(u(0,t)\) is a nonzero scalar multiple of
\(w\).  If \(s+t\ne0\), then
\[
 u(s+t,r)=2(s+t)w.
\]
These two decidable cases give the asserted projective parametrization
without extracting a square root.
\hfill\(\blacksquare\)

\medskip
\noindent\textbf{Theorem (Rational parabolic normal form).}
Fix a nonzero rational null vector \(u\), and let
\(\rho_u(\lambda)\) be the Minkowski rotation represented by
\((1:\lambda u)\).  Then
\[
\begin{split}
 \rho_u(0)&=I,\\
 \rho_u(\lambda+\mu)&=\rho_u(\lambda)\rho_u(\mu),\\
 \rho_u(-\lambda)&=\rho_u(\lambda)^{-1},
\end{split}
\]
and the map \(\lambda\mapsto\rho_u(\lambda)\) is injective.  Every member
fixes \(u\), has characteristic polynomial \((X-1)^3\), and has zero
handedness on every perpendicular transversal.

Conversely, every nonidentity member \(\rho\in SO(M,\mathbb Q)\) with
characteristic polynomial \((X-1)^3\) has the form
\[
 \rho=\rho_{u(\alpha,\beta)}(\lambda)
\]
for explicitly constructed \(\alpha,\beta,\lambda\in\mathbb Q\), with
\((\alpha,\beta)\ne(0,0)\) and \(\lambda\ne0\).  Once the displayed nonzero
axis vector is fixed, its parameter \(\lambda\) is unique.

\smallskip
\noindent\textbf{Proof.}
Since \(Q_M(u)=0\) and \(C_M(u,u)=0\), direct multiplication gives
\[
 (1:\lambda u)\otimes_M(1:\mu u)=(1:(\lambda+\mu)u),
 \qquad \overline{(1:\lambda u)}=(1:-\lambda u).
\]
This proves the three group laws.  If two resulting rotations agree,
projective faithfulness makes their representatives scalar multiples.
Their scalar coordinates are both one, so the scale is one; since \(u\ne0\),
their parameters agree.  The conjugation formula evaluated at \(u\) proves
that \(u\) is fixed.  Direct trace, principal-trace, and determinant
calculations give \((X-1)^3\), while expansion of the handedness determinant
gives zero.

For the converse, evaluate the supplied characteristic-polynomial identity
at \(X=1\) and \(X=-1\).  This gives trace three and hence half spread zero.
Choose the projective Clifford preimage \(P=(x:v)\) from the first theorem.
Its half-spread formula is
\[
 0=\frac{-Q_M(v)}{N_M(P)},
\]
so \(Q_M(v)=0\).  Non-nullness of \(P\) then forces \(x\ne0\), while
nonidentity forces \(v\ne0\).  Rescaling by \(x^{-1}\) gives
\((1:x^{-1}v)\).  The null-conic theorem supplies
\(\alpha,\beta\) and a nonzero scale relating \(x^{-1}v\) to
\(u(\alpha,\beta)\); absorbing that scale into a nonzero \(\lambda\) gives
the displayed normal form.  Uniqueness for a fixed axis is the injectivity
already proved.
\hfill\(\blacksquare\)

\medskip
\noindent\textbf{Corollary (Conjugacy of rational parabolics).}
Any two nonidentity rational Minkowski rotations with characteristic
polynomial \((X-1)^3\) are conjugate by a rational member of
\(SO(M,\mathbb Q)\).

\smallskip
\noindent\textbf{Proof.}
Apply the normal-form theorem to the two rotations.  For each parametrized
null vector, an explicit invertible \(2\)-by-\(2\) change of basis carries
its rank-one nilpotent matrix to
\[
 \begin{pmatrix}0&1\\0&0\end{pmatrix}.
\]
In these two bases the parabolics are upper unipotent.  The diagonal matrix
\(\operatorname{diag}(\lambda,\mu)\) intertwines the two nonzero parameters
\(\lambda\) and \(\mu\).  Multiplying the two basis changes and this diagonal
matrix gives an invertible rational transporter.  Lifting it through
\(\Phi^{-1}\) and taking its normalized Clifford action produces the required
rational special-orthogonal conjugator.
\hfill\(\blacksquare\)

\subsection{Problem 72: projective point--plane duality}

Problem 72 asks for a projective duality carrying metrical statements about
points to statements about planes, and for its application to the rotation
theorems of the paper.  Here a plane always means a plane through the origin.

\medskip
\noindent\textbf{Lemma (Representative point--plane polarity).}
Represent a projective point by a nonzero row vector \(x\), and a projective
plane by a nonzero row normal \(n\), in each case modulo multiplication by a
nonzero scalar.  The assignments
\[
 [x]\longmapsto x^\perp,\qquad
 n^\perp\longmapsto[n]
\]
are well defined and inverse to one another.  A point \([x]\) is incident
with \(n^\perp\) exactly when
\[
 x\mathbin{\cdot}n=0,
\]
and this condition is independent of both displayed representatives.

\smallskip
\noindent\textbf{Proof.}
The polar plane of \(x\) is defined by the equation
\(z\mathbin{\cdot}x=0\), so replacing \(x\) by
\(\lambda x\), with \(\lambda\ne0\), does not change that plane.
Conversely, replacing a normal \(n\) by \(\mu n\), with \(\mu\ne0\), does
not change either its plane or its projective point.  Both maps retain the
same three displayed coordinates; their two composites are therefore
literally the identity on representatives.  Finally,
\[
 (\lambda x)\mathbin{\cdot}(\mu n)
   =\lambda\mu(x\mathbin{\cdot}n),
\]
and \(\lambda\mu\ne0\), proving the representative independence of
incidence. \hfill\(\blacksquare\)

\medskip
\noindent\textbf{Theorem (Equivariance and incidence).}
Let \(\rho\in SO(3,F)\).  In the row-vector convention define the actions on
point and plane representatives by
\[
 x\longmapsto x\rho,\qquad n\longmapsto n\rho.
\]
These actions preserve nonzero representatives and projective equivalence,
satisfy the identity, product, and inverse laws, and commute with the
polarity:
\[
 (x\rho)^\perp=x^\perp\rho.
\]
Simultaneous rotation preserves incidence.

\smallskip
\noindent\textbf{Proof.}
Since \(\rho\rho^T=I\), multiplication by \(\rho^T\) recovers \(x\) from
\(x\rho\); hence a nonzero vector cannot be sent to zero.  Linearity gives
\((\lambda x)\rho=\lambda(x\rho)\), so the action respects projective
equivalence.  Associativity of matrix multiplication and
\(\rho^{-1}=\rho^T\) give the action laws.  Finally orthogonality gives
\[
 (x\rho)\mathbin{\cdot}(n\rho)
 =x\rho\rho^Tn^T
 =x\mathbin{\cdot}n.
\]
Thus the image of the plane \(x^\perp\) has normal \(x\rho\), which proves
equivariance, and the same identity proves incidence invariance.
\hfill\(\blacksquare\)

\medskip
\noindent\textbf{Theorem (Projective quadrance is dual-plane spread).}
Let \([x]\) and \([y]\) be projective points with
\(Q(x)\ne0\) and \(Q(y)\ne0\).  Then
\[
 q([x],[y])
 =1-\frac{(x\mathbin{\cdot}y)^2}{Q(x)Q(y)}
 =s(x^\perp,y^\perp).
\]
The value is independent of the chosen point or plane representatives and is
unchanged when both are acted on by the same rotation.

\smallskip
\noindent\textbf{Proof.}
By definition, the spread of two planes is the spread of their normal
vectors.  The polarity retains the coordinates \(x\) and \(y\), so its spread
is exactly the middle expression.  Rescaling \(x\) by \(\lambda\ne0\)
multiplies the numerator and the factor \(Q(x)\) in the denominator by
\(\lambda^2\), and these factors cancel; the calculation for \(y\) is the
same.  Rotation invariance follows from preservation of dot products and
quadrances.  The hypotheses \(Q(x),Q(y)\ne0\) are essential: over finite or
split fields, a nonzero representative may be null, in which case the
displayed quotient is not defined. \hfill\(\blacksquare\)

\medskip
\noindent\textbf{Theorem (Constructive translation of rotation theorems).}
Fix a rotation \(\rho\).  For every relation
\(\mathcal R(x,x\rho)\) on a nonzero vector and its image, the statements
\[
 \mathcal R(x,x\rho)\text{ for every projective point }[x]
\]
and
\[
 \mathcal R(n,n\rho)\text{ for every projective plane }n^\perp
\]
are equivalent.  More generally, for every four-place relation
\(\mathcal S(x,y,x\rho,y\rho)\), its universal point statement is equivalent
to the corresponding universal statement about two plane normals.

\smallskip
\noindent\textbf{Proof.}
Given the point statement, apply it to the projective point \([n]\) dual to
an arbitrary plane \(n^\perp\).  The rotation-equivariance theorem identifies
its image with the normal \(n\rho\) of the rotated plane.  This proves the
plane statement.  The reverse implication applies the plane statement to
the polar plane \(x^\perp\).  Since both polarity round trips are exact, the
original point statement is recovered.  Applying the same argument to two
representatives proves the four-place version. \hfill\(\blacksquare\)

\medskip
\noindent\textbf{Corollary (Plane forms of the principal rotation
theorems).}
The preceding translation gives the following dual statements.
\begin{enumerate}
\item A half turn fixes the plane whose non-null normal is its axis.  It also
fixes projectively every plane whose normal is perpendicular to that axis,
because such a normal is negated.
\item If two non-null independent vectors define the reflection product of
Theorem 19, the plane normal to their cross-product axis is fixed.
\item A defined cap fixes the plane normal to its nonzero anchor.  Away from
the half-turn chart, the recovered nonzero anchor of Theorem 35 is likewise
the normal of a fixed plane.
\item Let \(\rho\) have half spread \(p\).  Every non-null plane normal
perpendicular to the canonical rotation axis satisfies
\[
 s(n^\perp,n^\perp\rho)=4p(1-p).
\]
In particular, the value is independent of the supplied admissible
perpendicular normal.
\end{enumerate}

\smallskip
\noindent\textbf{Proof.}
For the first three clauses, use the corresponding vector equalities from
Theorems 17, 19, 24, and 35 as equalities of plane normals; a sign change is
a nonzero projective rescaling.  For the final clause, apply the
quadrance--spread correspondence to Theorem 31 and use equivariance of the
polarity.  The second spread polynomial in that theorem is
\(4p(1-p)\). \hfill\(\blacksquare\)

\section{Extent and boundary of the formal verification}

The formal verification covers the theorem-level finite algebraic content of
the paper, including its rational, projective, and finite-field
specializations.  It also covers the additional results stated in this note:
the corrected unital classification prompted by Problem 68, the finite-field
counts and classifications of Problem 69, the finite-query results for
Problem 70, the determinant-corrected ternary-form theory and rational
Minkowski classification of Problem 71, and the point--plane duality of
Problem 72.  Historical attribution, priority, novelty, and the unresolved
bibliographical placeholders in the source are editorial matters rather than
formal mathematical claims, and have not been verified.

Several distinctions are important.  The statements of Theorems 31, 40, 46,
and 50 are retained, but their printed proofs require the repairs given in
this note.  Corollary 49 requires the two additional input-cap domain
hypotheses displayed above, unless it is restricted to an ordered field where
they hold automatically.  Conjecture 52 is false, so Problem 68 as printed
has a negative answer; the unital homogeneous classification proved in
Section 8.2 is a corrected replacement, not a verification of the original
conjecture.
Likewise, the general-form results of Problem 71 require the determinant
factors displayed in Section 9.3.  These changes repair defective statements
or formulas; they are not concessions made in order to obtain a formal proof.

Other apparent changes are only changes of presentation.  Projective points,
planes, proportions, square classes, and cosets are represented by displayed
finite data together with explicit nonzero scale or equivalence witnesses.
The resulting statements are equivalent to the corresponding projective
ones; no conclusion is weakened by avoiding quotient objects.  Lemma 59 is
proved with its literal conclusion for every supplied finite number of
variables by the degree-plus-one argument of Section 6, and Theorem 58 is
proved independently by its direct coordinate identities.

The only source conclusions not repeated literally are those whose wording
requires a completed infinite collection or topology.  Their verified finite
content is as follows.
\begin{itemize}
\item For every supplied positive integer \(N\), the construction gives
\(N\) pairwise inequivalent positive rational square classes, each with an
explicit rotation preimage.  This replaces the infinitude sentence of
Theorem 62.
\item For every supplied positive integer \(N\), it gives \(N\) pairwise
kernel-separated rational rotations.  This replaces the infinite-index
sentence of Corollary 63 and its repetition in Section 13.4.
\item For every supplied rational rotation and every positive rational
squared-error bound, it constructs a unit-representable rational rotation
strictly within that bound.  This is the precise finite approximation result
proved in response to the density clause of Problem 70; no topological space,
closure, or density theorem is asserted.
\end{itemize}
Thus arbitrary finite size and arbitrary finite precision are retained, but
no completed carrier of all square classes, abstract quotient index, or
topological closure is formed.

Problems 69, 71, and 72 are solved in the finite senses stated in Sections
9.1, 9.3, and 9.4.  Problem 70 is only partly answered.  For every supplied
finite generator list and every supplied finite signed word, the rotation,
its \(\chi\)-value, kernel membership, and kernel coset are decided exactly,
and the complete finite \(\chi\)-image is constructed.  The finite-order
possibilities are also classified, subject in order two to the quadrance
class of the supplied rational axis.  What remains unproved is the internal
classification of the generally non-finite generated subgroups---their
isomorphism types, relations, or presentations---together with any broader
interpretation of ``visible through \(\chi\)'', a topological closure theorem,
and the global three-square classification needed to list all order-two axis
classes.  These are exactly the limitations recorded at the end of Section
9.2; no other finite algebraic result in this note is knowingly left
unverified.

\clearpage

\end{document}
